% Generated draft: requires mathematical and source review.
\documentclass[11pt]{article}
\usepackage[margin=25mm]{geometry}
\usepackage{fontspec}
\setmainfont{Latin Modern Roman}
\usepackage{amsmath,amssymb,mathtools,unicode-math}
\setmathfont{Latin Modern Math}
\usepackage{enumitem,needspace,xurl,hyperref}
\setlength{\parindent}{0pt}
\setlength{\parskip}{5pt}
\setlength{\emergencystretch}{3em}
\newcommand{\N}{\mathbb{N}}
\newcommand{\Z}{\mathbb{Z}}
\newcommand{\Q}{\mathbb{Q}}
\newcommand{\R}{\mathbb{R}}
\newcommand{\C}{\mathbb{C}}
\newcommand{\F}{\mathbb{F}}
\newcommand{\A}{\mathbb{A}}
\newcommand{\PP}{\mathbb{P}}
\newcommand{\E}{\mathbb{E}}
\newcommand{\eps}{\varepsilon}
\newcommand{\dd}{\,\mathrm{d}}
\begin{document}
\section*{UnsolvedMath ID: 10400231}
\textbf{Source problem code:} catalogue entry
\subsection{Definitions and mathematical statement}
Let \(H\) be a finitely generated abelian group with \(\dim_\Q(H\otimes\Q)=1\), let \(T=\operatorname{Tors}H\), and fix a generator \(t\) of \(H/T\cong\Z\). Then \(\Z[H/T]=\Z[t^{\pm1}]\). Let \(\Delta(t)\) be an element of this Laurent polynomial ring.

V. Turaev's conjecture asserts that the pair \((H,\Delta)\) is realizable by a closed connected oriented 3-manifold \(M\), with an isomorphism \(H_1(M;\Z)\cong H\) and matching Alexander polynomial, if and only if
\[
 \exists k\in2\Z:\quad\Delta(t)=t^k\Delta(t^{-1}),\qquad
 \Delta(1)=\pm|T|.
\]
Here \(|T|\) is the finite order of the torsion subgroup, equal to one when it is trivial. The symmetry sign is the positive sign printed in the source, and the parity restriction on \(k\) is retained.

In the usual rank-one Alexander convention, the map to \(H_1(M)/\operatorname{Tors}\cong\Z\) defines an infinite cyclic cover; \(\Delta_M\) is the order of its Alexander module, up to a unit \(\pm t^j\). Realizability is therefore understood with this unit ambiguity and the indicated free-generator identification. Multiplying by such a unit preserves evenness of the symmetry exponent and changes \(\Delta(1)\) only by sign. The source's exact Alexander normalization should still be checked.

The target is simultaneous realization of the prescribed abelian group, including all its torsion, and the prescribed polynomial. It is not merely realization of an Alexander polynomial with a specified rank, and it includes both necessity and sufficiency of the displayed conditions.
\subsection{Short English statement}
A prescribed rank-one abelian homology group and Alexander polynomial should be realizable by a closed oriented 3-manifold exactly when the polynomial has even-exponent reciprocal symmetry and evaluation at one equal, up to sign, to the torsion-group order.
\subsection{Sources}
\begin{itemize}
\item Local UnsolvedMath registry, permanent ID 10400231, source code catalogue entry; statement attribution V. Turaev; proposer field and proposal year absent; updated 2026-08-06.
\item Ohtsuki, \emph{Problems on invariants of knots and 3-manifolds} (2002), Conjecture 12.26, PDF page 170: \url{https://msp.org/gtm/2002/04/gtm-2002-04-024s.pdf}. Registry access date 2026-07-29; PDF font-structure extraction. Provenance distinguishes 231 numbered bold headings. Present primary-literature status and redistribution terms require review.
\item The supplied summary cites V. Turaev, ``Torsions of 3-manifolds'' (2002), arXiv:math/0211084, for related torsion-realization questions. Its reported partial rank-two/rank-three augmentation-one realizations concern a different setting and are not adopted as a solution of this rank-one pair-realization conjecture. The citation and status report were not independently verified.
\end{itemize}
\end{document}
